% ====================
% 4. System or study design
% ====================
\section{System or Study Design}
\label{sec:methodology}

\subsection{Scope Representation and Conflict Resolution}
\label{sec:scope-representation}

Each instruction is represented by three fields:
\emph{scope}, \emph{key}, and \emph{condition}.
The scope determines applicability, the optional key identifies
a project or task, and the optional condition specifies an
additional constraint.

\begin{table}[t]
\centering
\small
\caption{Scope representation.}
\label{tab:scopes}
\begin{tabularx}{\linewidth}{lXl}
\toprule
Scope & Applicability & Example \\
\midrule
Standing & General preference across tasks &
Use Polars \\
Project & Restricted to a named project &
Marlin: pandas \\
Single-task & Restricted to one task &
This export: CSV \\
\bottomrule
\end{tabularx}
\end{table}

The reference decision policy first excludes memories whose
project, task, or condition does not match the current request.
Among applicable memories concerning the same decision, scope
precedence is
\[
\text{single-task} \succ \text{project}
\succ \text{standing}.
\]
Within one scope, an applicable conditional instruction takes
precedence over an unconditional instruction; recency resolves
remaining conflicts at equal specificity. Unrelated instructions
do not compete, and the benchmark excludes unresolved ties.

These rules define the reference actions and are communicated
to the agent. The scope layer does not deterministically
select the winning memory.

\subsection{Scope-Tagging Architecture}
\label{sec:architecture}

We augment Mem0~\cite{chhikara2025mem0} and
Supermemory~\cite{supermemoryresearch} with an LLM-based
scope tagger. Figure~\ref{fig:architecture} compares the
two integration strategies.


\begin{figure}[t]
\centering
\begin{tikzpicture}[yscale=0.62,
    >=Latex,
    node distance=0.55cm,
    box/.style={
        draw=gray!60,
        rounded corners=2pt,
        minimum width=3.9cm,
        minimum height=0.5cm,
        align=center,
        font=\footnotesize
    },
    tag/.style={
        box,
        draw=blue!70,
        fill=blue!8,
        font=\footnotesize\bfseries
    },
    arrow/.style={->,thick,gray!70}
]

% Column headings
\node[font=\bfseries] at (0,0) {Write-time};
\node[font=\bfseries] at (5.5,0) {Read-time};

% Write-time pipeline
\node[box] (w1) at (0,-1) {Conversation};
\node[tag] (w2) at (0,-2) {Extract and tag};
\node[box] (w3) at (0,-3) {Memory store};
\node[box] (w4) at (0,-4) {Retrieve};
\node[box] (w5) at (0,-5) {Agent + resolution rules};

% Read-time pipeline
\node[box] (r1) at (5.5,-1) {Conversation};
\node[box] (r2) at (5.5,-2) {Memory store};
\node[box] (r3) at (5.5,-3) {Retrieve};
\node[tag] (r4) at (5.5,-4) {Tag retrieved text};
\node[box] (r5) at (5.5,-5) {Agent + resolution rules};


% Write-time arrows
\draw[arrow] (w1.south) -- (w2.north);
\draw[arrow] (w2.south) -- (w3.north);
\draw[arrow] (w3.south) -- (w4.north);
\draw[arrow] (w4.south) -- (w5.north);

% Read-time arrows
\draw[arrow] (r1.south) -- (r2.north);
\draw[arrow] (r2.south) -- (r3.north);
\draw[arrow] (r3.south) -- (r4.north);
\draw[arrow] (r4.south) -- (r5.north);


\end{tikzpicture}
\caption{Scope-tagging integration strategies.
Write-time tagging occurs before storage, whereas
read-time tagging occurs after retrieval.}
\label{fig:architecture}
\end{figure}


Write-time tagging extracts individual instructions from the
full conversation and stores them with scope metadata. It
can preserve scope cues early but may discard surrounding
conversational context. Read-time tagging operates on the
memory system's retrieved text, preserving native storage
behaviour but relying on that text to retain sufficient
contextual information.

Both strategies supply tagged memories and conflict-resolution
rules to the agent. Structured task identifiers are withheld
from the taggers, so task continuity must be inferred from
natural-language context.

\subsection{Benchmark Construction}
\label{sec:benchmark}

\subsubsection{Episode Structure and Mirrored-Pair Design}

We construct 20 paired scenarios containing 40 episodes.
Each pair includes two or three earlier conversation sessions,
reference memories, and two probes sharing the same history.
Probe A tests an applicable scoped instruction; Probe B
changes the project, task, or condition so that the instruction
no longer applies. Each episode consists of one probe and
its shared history, with an expected action for evaluation.

Each pair also includes an unrelated standing preference and
a retention check. Pair accuracy requires both probes to be
correct, discouraging fixed choices that ignore context.

Appendix~\ref{app:po01-example} (Figure~\ref{fig:mirrored-pair}) walks through pair \texttt{po-01}.

\subsubsection{Conflict Categories and Dataset Generation}

The benchmark covers five conflict categories with four
pairs each, testing project exceptions, temporary task
exceptions, conditional applicability, replacement of older
instructions within one scope, and precedence of more
specific instructions within one scope.

The 20 pairs contain 66 distinct reference memories,
including competing and unrelated instructions.
Nineteen pairs were initially drafted with AI assistance
and subsequently reviewed and revised. The suite varies
domains, wording, chronology, and instruction positions;
some standing preferences deliberately appear after
narrower instructions.

\subsubsection{Dataset Validation and Gold Labels}

Automated checks validate the JSON schema, verbatim
memory quotations, paired-probe similarity, expected
decisions, and diversity. Only session and probe text
reach the systems; structured gold labels and expected
actions are withheld. A leakage test checks this
separation.

The 66 reference memories were independently annotated
by two human annotators for scope, key, and condition.
The annotations were reviewed and reconciled to establish
the final gold labels. After reconciliation, scope and key
labels matched the original drafts, while four
condition-expression differences were reviewed and resolved. 

